% paper/sections/11b_force_balance.tex
%
% §11.1  力学的バランス検証（Force Balance Verification）
%   §11.1.1 静水圧均衡（重力 vs 圧力勾配）
%   §11.1.2 表面張力均衡（Laplace 圧 vs 寄生流れ）
% ═══════════════════════════════════════════════════════════════════════════════
\subsection{力学的バランス検証}
\label{sec:force_balance}
% ═══════════════════════════════════════════════════════════════════════════════

NS ソルバの最も基本的な物理的整合性は，
\textbf{静止流体における力の釣合が数値的に維持されること}である．
本節では 2 種類の均衡テストにより検証する．

% ─────────────────────────────────────────────────────────────────────────────
\subsubsection{静水圧均衡}
\label{sec:hydrostatic}
% ─────────────────────────────────────────────────────────────────────────────

%% 実験結果は exp11_1_hydrostatic.py の実行後に挿入

\textbf{設定：}
計算領域 $[0,1]^2$，単相流（$\psi = 1$），壁面境界条件（上下左右），
重力 $\bm{g} = (0, -1)^\top$．初期条件 $\bu = \bm{0}$，$p = 0$．
格子解像度 $N = 32, 64, 128$．

\textbf{理論解：} 圧力は静水圧分布 $p(y) = \rho g (1-y)$ に収束し，
速度は恒等的にゼロ（$\bu \equiv \bm{0}$）であるべきである．

\textbf{合格基準：} $\|\bu\|_\infty < 10^{-12}$（計算機精度水準の寄生速度）．

%% TODO: exp11_1 結果テーブルをここに挿入

% ─────────────────────────────────────────────────────────────────────────────
\subsubsection{表面張力均衡と寄生流れ}
\label{sec:bf_static_droplet}
% ─────────────────────────────────────────────────────────────────────────────

%% 実験結果は exp11_2_bf_droplet.py の実行後に挿入

\textbf{設定：}
計算領域 $[0,1]^2$ に半径 $R = 0.25$ の円形液滴（$\rho_l/\rho_g = 1000$，$\mathit{We} = 1$）を配置する．
重力なし，壁面境界条件．初期条件 $\bu = \bm{0}$，初期圧力は Laplace 圧 $p^0 = \sigma\kappa_0 = \sigma/R$
から設定する．格子解像度 $N = 32, 64, 128$．

\textbf{理論解：}
Young--Laplace 条件 $\Delta p = \sigma/R = 4.0$（$R=0.25$, $\sigma=1$）．
速度は恒等的にゼロであるべきだが，離散化誤差により寄生流れ（spurious currents）が発生する．

\textbf{合格基準：}
\begin{itemize}[noitemsep]
  \item Laplace 圧相対誤差 $|\Delta p - \sigma/R|/(\sigma/R) < 5\%$
  \item 寄生流れ振幅の格子収束次数 $\geq 2$（Balanced-Force 条件下で $\Ord{h^2}$）
\end{itemize}

%% TODO: exp11_2 結果テーブルをここに挿入
